%% guide-p2a.tex　ガイドブック　保護者・先生のための本（大人用、vol=0）第2〜10章
%% （第16回編集委員会、組 G-B、頭文字 kgb）。ルビなし。大人向けの漢字の使い方（10月30日 10:15 依頼者）。
%% 材料：tools/guide-shidou.md（RG-1）・tools/guide-series.md（RG-2）・tools/guide-taiou.md（RG-3）。
%% 数・事実は材料のとおり。材料で「未確認」とされたことは書かない（または断定しない）。
%% 部品（\kgbsec・\kgbsub・\kgbvol・表の型 kgb など）は guide-a1.tex（第1章、先に読む）の先頭で定義している。
%% 出典は guide-p2a-src.tex（親ファイルで付録 guide-p2b のあとに読む）。
\ifdefined\kgbsec\else\errmessage{guide-p2a.tex の前に guide-a1.tex を読んでください（部品の定義がそこにあります）}\fi

%% ====================================================================
\kgbsec{2}{編集方針と書き方の原則}\label{kgb-ch2}

この章では、編集委員会が第1回の委員会で決めた編集方針と、全巻で守っている書き方の原則を紹介します。シリーズの目的・巻の並び・紙面の部品は、第1章にまとめました。

\kgbsub{2.1}{編集方針の10項目}

次の10項目が編集方針です。どの巻も、この方針に沿って書かれています。

\par\addvspace{4pt}
\kgbcap{編集方針の10項目と、紙面での表れ}{}
\begin{bktable}[long, theme=kgb]{hline{Z}={0pt,white}, vline{1-Z}={Z}{0pt,white}, row{Z}={abovesep=3pt, belowsep=1pt}, hspan=minimal, rowhead=1, colspec={Q[c,m,wd=6mm] Q[l,m,wd=62mm] Q[l,m,wd=58mm]}, rowsep=2pt, colsep=4pt, cells={font=\small}, row{1}={font=\sffamily\bfseries\small}}
 & 編集方針 & 紙面では \\
1 & 量＝数値×単位から出発する。「5」ではなく「5 m」 & 2-1・2-10。各巻の「3つのやくそく」の1つめ「数にはいつも単位」 \\
2 & 具体→表→グラフ→式の順。式は「表とグラフを短く書く方法」として最後に出す & 3-0 の「はかる→表→点→線→式」の5つの段 \\
3 & 大学まで続く種を早くまく（ひき算＝変化量、かけ算＝面積、九九の表＝行列）。7本の路線でつなぐ & 1-5「かわった量」、1-13 九九の表（行列の種）、7本の路線（1.3） \\
4 & 微分方程式は「次の瞬間を決めるルール」。式で解く前に表計算で解く & 第6巻を全巻の乗りかえ駅にした（6-2 オイラー法） \\
5 & 計算の量より意味。九九と筆算は第1巻で手でできるまで練習し、第2巻からは電卓・表計算・プログラムに任せる & 電卓は 2-1 から、表計算は 2-18 から、プログラムは 2-19 から使う。今の版の第1巻の筆算は足し算・引き算だけで、掛け算・割り算の筆算はない（7.2） \\
6 & どの章も実験や現象とセット。大学の章も、家や教室で確かめられる現象から入る & 「やってみよう」「物理・工学ではこう使う」 \\
7 & 物理や工学で使わない単元は省くか軽くする（統計は戻した） & 7章にまとめた \\
8 & 大学の内容も中学生の言葉で。「目でわかる→手でわかる→式でわかる」の3段で、前の2段は★★★以下の道具だけ & 6-10〜6-14、7-7〜7-10、8-1〜8-13、9-1〜9-11（1.5） \\
9 & 専門書の書き方に慣れさせる。新しい記号は1章に2つまでで、読み方・意味・英語を記号カードで示す & 記号カード、★★★★以上の章の「専門書をのぞいてみよう」 \\
10 & 学年ではなく、★と道しるべで進む & 章の帯の★、帯の下の「道しるべ」（前に読む章・この先の章） \\*
\SetCell[c=3]{c} \kgbcapline & & \\
\end{bktable}
\par\addvspace{4pt}

\kgbsub{2.2}{書き方の原則}

新しい考えは、いつも次の順で書いています。新しい言葉は体験の後で「このことを〇〇といいます」と1回だけ紹介し、名前を最後にします。

\begin{bkfigure}
\begin{tikzpicture}[x=1mm, y=1mm]
  \foreach \kgbi/\kgbt/\kgbf [count=\kgbn] in {0/場面/white,1/ふしぎ・\\こまったこと/white,2/手を動かす/white,3/気付く/white,4/言葉・\\記号・式/bkVol0!10} {
    \node[kgb box, minimum width=22mm, minimum height=13mm, text width=20mm, fill=\kgbf] (k\kgbn) at ({\kgbi*28.5+12},0) {\kgbt};}
  \foreach \kgbn [evaluate=\kgbn as \kgbm using int(\kgbn+1)] in {1,...,4} {\draw[kgb arr] (k\kgbn) -- (k\kgbm);}
  \node[kgb lab, anchor=north, text width=24mm, align=center] at (69,-9) {測る・数える・\\表にする・図をかく};
  \node[kgb lab, anchor=north, text width=22mm, align=center] at (126,-9) {名前は最後};
\end{tikzpicture}
\bkfigcaption{新しい考えを紹介する順番（名前は最後）}
\end{bkfigure}

このほか、次のことを全巻で守っています。

\kgbitem{公式は作ってから使う。表や図で、いくつかの数で同じことが起きるのを確かめてから、決まりにします。}
\kgbitem{数字と文章だけで説明しない。新しい考えには必ず図を付けます。例題は、考え方と途中の式を1行ずつ全部書きます。}
\kgbitem{式が出てくるたびに「なぜその式になるのか」、数が出てくるたびに「なぜその数なのか」を書きます。初めての考え方は2通り以上で説明します。}
\kgbitem{例題で見せていない形を、「じぶんでとこう」や章末問題に出しません。}
\kgbitem{数学として後で困る言い方をしません（「かけ算すると大きくなる」と言い切らない、「0でわる」ことはできないと書く、など）。たとえ話には「せいかくに言うと」を添えます。}
\kgbitem{工学につなげます。本物の機械やしくみで、その章の数学がどう働いているかを見せます。}

%% ====================================================================
\kgbsec{3}{学習指導要領との関係のあらまし}\label{kgb-ch3}

\kgbsub{3.1}{現行の学習指導要領}

この本で照らし合わせたのは、次の現行の学習指導要領（算数・数学）です。授業時数は学校教育法施行規則の別表、高校の単位数は指導要領の標準単位数です。

\par\addvspace{4pt}
\kgbcap{現行の学習指導要領（算数・数学）}{}
\begin{bktable}[long, theme=kgb]{hline{Z}={0pt,white}, vline{1-Z}={Z}{0pt,white}, row{Z}={abovesep=3pt, belowsep=1pt}, hspan=minimal, rowhead=1, colspec={Q[l,m,wd=17mm] Q[l,m,wd=27mm] Q[l,m,wd=36mm] Q[l,m,wd=46mm]}, rowsep=2pt, colsep=4pt, cells={font=\small}, row{1}={font=\sffamily\bfseries\small}}
学校 & 告示と実施 & 授業時数・単位数 & 領域・科目 \\
小学校（算数） & 平成29年告示、令和2年度から全面実施 & 計1011時間（1年136、2〜6年は各175。1単位時間45分） & A 数と計算／B 図形／C 測定（1〜3年）・C 変化と関係（4〜6年）／D データの活用 \\
中学校（数学） & 平成29年告示、令和3年度から全面実施 & 計385時間（1年140、2年105、3年140。1単位時間50分） & A 数と式／B 図形／C 関数／D データの活用 \\
高等学校（数学） & 平成30年告示、令和4年度から学年進行で実施 & 数学Ⅰ 3単位（必履修）、数学Ⅱ 4、数学Ⅲ 3、数学A 2、数学B 2、数学C 2 & 数学Ⅰ・Ⅱ・Ⅲ・A・B・C（「数学活用」をなくし、数学Cを新しく置いた） \\*
\SetCell[c=4]{c} \kgbcapline & & & \\
\end{bktable}
\par\addvspace{4pt}

高校の6科目は、次の図のように組み立てられています。数学A・B・Cは、それぞれ3つの項目から適宜選んで学ぶ科目です。

\begin{bkfigure}
\begin{tikzpicture}[x=1mm, y=1mm]
  \tikzset{kgb sub/.style={kgb gbox, text width=38mm, anchor=north}}
  \node[kgb sub, draw=bkVol0, line width=1pt, minimum height=31mm] (s1) at (21,0) {{\sffamily\bfseries 数学Ⅰ}（必履修）\\(1) 数と式\\(2) 図形と計量\\(3) 二次関数\\(4) データの分析};
  \node[kgb sub, minimum height=31mm] (s2) at (70,0) {{\sffamily\bfseries 数学Ⅱ}\\(1) いろいろな式\\(2) 図形と方程式\\(3) 指数関数・対数関数\\(4) 三角関数\\(5) 微分・積分の考え};
  \node[kgb sub, minimum height=31mm] (s3) at (119,0) {{\sffamily\bfseries 数学Ⅲ}\\(1) 極限\\(2) 微分法\\(3) 積分法};
  \draw[kgb arr] (s1.east |- 0,-15.5) -- (s2.west |- 0,-15.5);
  \draw[kgb arr] (s2.east |- 0,-15.5) -- (s3.west |- 0,-15.5);
  \foreach \kgbx/\kgbn/\kgba/\kgbb/\kgbc/\kgbo in {21/数学A/図形の性質/場合の数と確率/数学と人間の活動/数学Ⅰと並行して、または後に,
    70/数学B/数列/統計的な推測/数学と社会生活/数学Ⅰの後が原則,
    119/数学C/ベクトル/平面上の曲線と複素数平面/数学的な表現の工夫/数学Ⅰの後が原則} {
    \node[kgb sub, minimum height=29mm] at (\kgbx,-45) {{\sffamily\bfseries \kgbn}\\(1) \kgba\\(2) \kgbb\\(3) \kgbc\\[2pt]{\sffamily\color{bkGray}\kgbo}};}
  \node[kgb lab, anchor=west, text=bkGray] at (2,-41) {数学A・B・C：それぞれ3つの項目から適宜選ぶ科目};
\end{tikzpicture}
\bkfigcaption{高校の数学の6科目と、学ぶ順序（矢印：数学Ⅰ→Ⅱ→Ⅲの順が原則）}
\end{bkfigure}

\kgbnote{※ どの学年でどの科目を取るかは、指導要領では決まっていません。一般的な例では、1年で数学Ⅰ・数学A、2年で数学Ⅱ・数学B、3年で数学Ⅲ・数学Cを取ります（数学Cの一部を2年に置く学校もあります）。この本で高校の「学年」を書くときは、この一般的な例によっています。}
\kgbnote{※ 令和7年度の大学入学共通テストでは、数学A・B・C のうち (3) の項目（数学と人間の活動・数学と社会生活・数学的な表現の工夫）は出題範囲外です。行列は、このうち数学Cの (3)「数学的な表現の工夫」に入っています（6.2を参照）。}

\kgbsub{3.2}{このシリーズの立場}

このシリーズは、学習指導要領の学年の並びに沿うのではなく、物理・工学で使う道筋（7本の路線）に沿って内容を並べています。そのため、指導要領の学年とは、次の3つの点で一致しない所があります。

\kgbitem{{\sffamily\bfseries 先に出る内容}　指導要領では上の学年の内容が、前の巻に出てきます。例：小4の折れ線グラフ・角・面積（第1巻 1-18・1-20・1-21）、中3の平方根（第2巻 2-8）、数学Ⅱの指数関数・対数（第3巻 3-11・3-12）。}
\kgbitem{{\sffamily\bfseries 後に出る内容}　指導要領では下の学年の内容が、後の巻に出てきます。例：小3〜小6の図形の面積・体積・合同・対称（第4巻 4-1・4-3・4-5）、小6の代表値（第3巻 3-13）。}
\kgbitem{{\sffamily\bfseries 分かれている内容}　指導要領の1つの単元が、いくつかの巻・章に分かれています（表\ref{kgb-t-tangen}）。}

\par\addvspace{4pt}
\kgbcap{1つの単元が、いくつかの章に分かれている例}{kgb-t-tangen}
\begin{bktable}[long, theme=kgb]{hline{Z}={0pt,white}, vline{1-Z}={Z}{0pt,white}, row{Z}={abovesep=3pt, belowsep=1pt}, hspan=minimal, rowhead=1, colspec={Q[l,m,wd=30mm] Q[l,m,wd=98mm]}, rowsep=1.5pt, colsep=4pt, cells={font=\small}, row{1}={font=\sffamily\bfseries\small}}
単元 & このシリーズで学ぶ章（学ぶ順） \\
面積 & 1-21（長方形）→ 4-5（平行四辺形・三角形・台形・円・柱・すい・球）。円周率は 2-9、おうぎ形は 4-2 \\
平均・データ & 2-6（測定値の平均）→ 3-13（度数分布・代表値・箱ひげ図）→ 3-14・3-15 \\
比例・関数 & 2-15（比例・反比例）→ 3-3（関数）→ 3-5（1次関数）→ 3-7（反比例のグラフ）→ 3-9（$y=ax^2$） \\
座標 & 3-2（平面）→ 4-13（空間）→ 4-15（極座標・円柱座標・球座標）。極座標は 4-2 でも紹介 \\
平方完成 & 3-8（式）・3-9（頂点）・3-10（解の公式）・3-16（最小二乗法） \\
指数 & 2-7（累乗・指数のきまり）→ 3-11（指数関数・分数の指数）→ 5-14（$e$） \\
三角比・三角関数 & 4-6 → 4-7 → 4-8 → 4-9 → 5-15 \\
数列にあたるもの & $\Sigma$ は 5-10、漸化式と等比数列・等比級数は 6-2（オイラー法）と 7-10 \\
確率 & 3-17（場合の数は確率に必要な分だけ） \\*
\SetCell[c=2]{c} \kgbcapline & \\
\end{bktable}
\par\addvspace{4pt}

\kgbsub{3.3}{用語と書き方の違い}

工学の習慣に合わせたり、子どもに分かりやすい言葉を選んだりしたため、学校の教科書と用語・書き方が違う所があります。

\par\addvspace{4pt}
\kgbcap{学校の教科書と、用語・書き方が違う所}{kgb-t-yougo}
\begin{bktable}[long, theme=kgb]{hline{Z}={0pt,white}, vline{1-Z}={Z}{0pt,white}, row{Z}={abovesep=3pt, belowsep=1pt}, hspan=minimal, rowhead=1, colspec={Q[l,m,wd=40mm] Q[l,m,wd=50mm] Q[l,m,wd=36mm]}, rowsep=1.5pt, colsep=4pt, cells={font=\small}, row{1}={font=\sffamily\bfseries\small}}
学校の教科書 & このシリーズ & 章 \\
虚数単位 $i$ & $j$（$i$ は電流に使うため） & 6-8 から \\
複素数平面 & 複素平面 & 6-9 \\
楕円 & だ円 & 4-8・4-15 ほか \\
概数 & がい数 & 2-6 \\
ドットプロット & 点を並べた図（名前は付けていない） & 2-6 \\
素因数分解 & 素数のかけ算に分ける（語はない） & 2-4 \\
最大公約数・最小公倍数 & 語はない（「公倍数で一番小さい数」など） & 2-4 \\
二項定理 & 語はない（パスカルの三角形） & 3-8・5-4（もっと知りたい） \\
二項分布・確率変数・確率分布 & 語はない（コインの表の枚数の分布） & 3-17 \\
漸化式 & 6-2 では語を使わず、7-10 で使う & 6-2・7-10 \\
変曲点 & 語はない（2回微分で山と谷を見分ける） & 5-7 \\
2進数の書き方 & $1011_{(2)}$ & 2-17 \\*
\SetCell[c=3]{c} \kgbcapline & & \\
\end{bktable}
\par\addvspace{4pt}

%% ====================================================================
\kgbsec{4}{学年と巻の対応}\label{kgb-ch4}

\kgbsub{4.1}{学年と巻の対応の図}

次の図は、指導要領の学年・科目の内容が、どの巻にあるかを示したものです。塗りつぶしたマスは、その学年・科目の内容（一部だけの場合も含みます）がその巻にあることを、太い枠は、その巻の学年の目安（4.2）を表します。高校の欄は科目ごとに分け、上に一般的な学年の例を添えました（学年の目安の枠も、この例で置いています）。一番右の「大学の内容」は、指導要領のどこにもなく、大学の工学部で学ぶ内容です。

\begin{bkfigure}
\begin{tikzpicture}[x=1mm, y=1mm]
  \kgbgridhead{1}
  \kgbgridrow{-5.0}{1}{1,2,3,4}{1}{3}
  \kgbgridrow{-13.2}{2}{2,3,4,5,6,7,8,9,10,11,12}{4}{7}
  \kgbgridrow{-21.4}{3}{4,6,7,8,9,10,11,12,13,14,15,16}{7}{9}
  \kgbgridrow{-29.6}{4}{3,4,5,6,7,8,9,10,11,12,15,16}{8}{13}
  \kgbgridrow{-37.8}{5}{11,12,13,14,16}{10}{13}
  \kgbgridrow{-46.0}{6}{12,13,14,15,16}{12}{15}
  \kgbgridrow{-54.2}{7}{13,14,15,16}{14}{16}
  \kgbgridrow{-62.4}{8}{14,15,16}{16}{16}
  \kgbgridrow{-70.6}{9}{12,14,16}{16}{16}
  \kgbgridlines{-75.2}
  % 凡例
  \fill[bkVol0!70, draw=bkInk!55, line width=0.4pt] (2,-82.7) rectangle ++(5.4,5.4);
  \node[kgb lab, anchor=west] at (9,-80) {内容がある（一部だけの場合も含む）};
  \draw[bkInk, line width=1.3pt, rounded corners=1.3mm] (78,-83.6) rectangle ++(13.8,7.2);
  \node[kgb lab, anchor=west] at (93.5,-80) {学年の目安};
\end{tikzpicture}
\bkfigcaption{学年・科目と巻の対応（マスは表\ref{kgb-t-gakunen} を基に作った）}\label{kgb-f-grid}
\end{bkfigure}

図を見ると、次のことが分かります。

\kgbitem{第1〜3巻は、ほぼ目安の学年の内容です。ただし第1巻には小4の内容（折れ線グラフ・角・面積）が、第2・3巻には中3〜高校の内容が入っています。}
\kgbitem{第4巻は、小3の図形から大学の内容（外積、円柱座標・球座標）までを含む、一番幅の広い巻です。第2巻は小2〜高校、第3巻は小4〜高校の内容を含みます。}
\kgbitem{第5巻の目安は「高校1年〜2年」ですが、数学Ⅲ（一般的な例では3年）の微分法・積分法も入っています。}
\kgbitem{第6巻からは、ほとんどが指導要領の外の内容です。高校の欄にあるますは、複素数（6-8・6-9）や、行列の章の中の等比数列・無限等比級数（7-10）などです。}

\kgbsub{4.2}{9巻の学年の目安}

各巻には、編集のときに決めた「学年の目安」があります。生徒用の紙面には学年を書かない決まりにしています（上の学年の子は恥ずかしく、下の学年の子は手を出さなくなるため）。学年の目安は、この本と、今後作る指導書にだけ書きます。各巻の学年の目安は、表\ref{kgd-t-vols}（1.2）の通りです。

\kgbsub{4.3}{学年・科目ごとの、主な巻と章}

図\ref{kgb-f-grid} のマスを、章まで示しました（章ごとの対応は付録A・付録B）。

\par\addvspace{4pt}
\kgbcap{学年・科目の内容が置かれている巻と、主な章}{kgb-t-gakunen}
\begin{bktable}[long, theme=kgb]{hline{Z}={0pt,white}, vline{1-Z}={Z}{0pt,white}, row{Z}={abovesep=3pt, belowsep=1pt}, hspan=minimal, colspec={Q[l,m,wd=12mm] Q[l,m,wd=116mm]}, rowhead=1, rowsep=1.5pt, colsep=4pt, cells={font=\footnotesize}, row{1}={font=\sffamily\bfseries\footnotesize}}
学年 & 置かれている巻（主な章） \\
小1 & 第1巻のみ \\
小2 & 第1巻（＞・＜だけ第2巻 2-2） \\
小3 & 第1巻（中心）、第2巻（□の式 2-11、k・t 2-10、万→億 2-7）、第4巻（二等辺三角形・正三角形 4-1、球 4-5） \\
小4 & 第1巻（折れ線グラフ 1-18、角 1-20、面積 1-21、余りの式 1-15）、第2巻（億・兆、概数、小数・分数、式、割合）、第3巻（ものの位置 3-2）、第4巻（平行・垂直 4-1、立体 4-5・4-13） \\
小5 & 第2巻（数と計算・速さ・割合・比例・平均・円周率）、第4巻（合同・多角形・面積・体積・角柱と円柱） \\
小6 & 第2巻（分数の乗除・文字・比・比例と反比例）、第3巻（代表値・度数分布・起こり得る場合）、第4巻（拡大図と縮図・対称・円の面積・体積） \\
中1 & 第2巻（正負の数・文字式・一次方程式・比例と反比例・素数）、第3巻（座標・関数・比例式・ヒストグラム・確率）、第4巻（平面図形・空間図形） \\
中2 & 第2巻（式の計算 2-12、等式の変形 2-14）、第3巻（連立方程式・一次関数・箱ひげ図・確率）、第4巻（平行線と角・合同・証明） \\
中3 & 第2巻（平方根・無理数・誤差・$a\times10^n$）、第3巻（展開・因数分解・二次方程式・$y=ax^2$・標本調査）、第4巻（相似・円周角・三平方の定理） \\
数Ⅰ & 第2巻（実数 2-9）、第3巻（数と式・二次関数・データの分析）、第4巻（図形と計量） \\
数A & 第2巻（2-4・2-17）、第3巻（3-17）、第4巻（4-1・4-2・4-7・4-11・4-13）、第5巻（$n!$ 5-17） \\
数Ⅱ & 第2巻（指数の拡張 2-7・2-12）、第3巻（3次の展開・恒等式・判別式・直線・指数と対数）、第4巻（距離・円の式・内分点・三角関数）、第5巻（微分・積分）、第6巻（複素数・解と係数の関係）、第9巻（分数式・3次方程式の一部） \\
数B & 第3巻（3-14・3-16・3-17）、第5巻（$\Sigma$ 5-10）、第6巻（漸化式の形 6-2）、第7巻（7-10） \\
数Ⅲ & 第3巻（合成関数・逆関数・$y=a/x$）、第5巻（極限・微分法・積分法）、第6巻（曲線の長さ：せいかくに言うと）、第7巻（$r^n$・無限等比級数）、第8巻（長さでたす線積分）、第9巻（無限級数としての $e^{At}$） \\
数C & 第3巻（双曲線の名 3-7、行列の種 3-6）、第4巻（ベクトル・極座標・媒介変数・だ円）、第6巻（複素平面 6-9）、第7巻（行列）、第8巻（だ円の式 8-1、媒介変数 8-7） \\*
\SetCell[c=2]{c} \kgbcapline & \\
\end{bktable}
\par\addvspace{4pt}

%% ====================================================================
\kgbsec{5}{巻ごとの対応}\label{kgb-ch5}

巻ごとに、指導要領のどの学年・科目に当たるかをまとめます（カードの帯は図\ref{kgb-f-grid} のその巻の行。章ごとの対応は付録B）。

\begin{kgbcard}{1}{かぞえる・はかる}
\kgbstrip{1}{1,2,3,4}{1}{3}
\kgbrow{学年の目安}{小学1年〜3年　{\sffamily ★の範囲} \kgbst{1}{1}　22章・400ページ}
\kgbrow{当たる範囲}{小1〜小3の数と計算（足し算・引き算、掛け算の意味と九九、割り算、小数と分数の入口）、測定（長さ・かさ・重さ・時刻と時間）、データ（表・棒グラフ）、図形（三角形・四角形・円）。小4の折れ線グラフ（1-18）・角（1-20）・面積（1-21）}
\kgbrow{目安とのずれ}{小4の内容が3章（★の基準「小4まで」の内）。小3の一部（そろばん、掛け算の筆算、二等辺三角形、km、□の式）は入っていないか、後の巻にある}
\kgbrow{指導要領の外}{引き算を「かわった量」（変化量）として見る見方（1-5）は、このシリーズ独自です}
\end{kgbcard}

\begin{kgbcard}{2}{数と量のことば}
\kgbstrip{2}{2,3,4,5,6,7,8,9,10,11,12}{4}{7}
\kgbrow{学年の目安}{小学4年〜中学1年　{\sffamily ★の範囲} \kgbst{1}{2}　20章・456ページ}
\kgbrow{当たる範囲}{小4〜小6の数と計算・変化と関係（割合・速さ・比・比例と反比例）・平均と概数、中1の数と式（正負の数・文字式・一次方程式）。中3の平方根・無理数（2-8・2-9）、誤差と $a\times10^n$（2-6・2-7）、数学Ⅰの実数（2-9）、数学Ⅱの0乗・負の指数（2-7・2-12）、数学Aの内容の取扱いで題材の例に挙がる二進法（2-17の2進数）}
\kgbrow{目安とのずれ}{中3〜高校の内容が5章。小4〜小6の図形は第4巻へ、データ（代表値・度数分布）は第3巻へ移した}
\kgbrow{指導要領の外}{量＝数値×単位、とびとびの数と切れ目のない数（2-1）、変化量 $\Delta$（2-3）、SI単位・組立単位・接頭語 μ（2-10）、16進数（2-17）、表計算（2-18）、プログラム（2-19）}
\end{kgbcard}

\begin{kgbcard}{3}{関数とデータ（きまりを見つける）}
\kgbstrip{3}{4,6,7,8,9,10,11,12,13,14,15,16}{7}{9}
\kgbrow{学年の目安}{中学1年〜3年　{\sffamily ★の範囲} \kgbst{1}{3}　18章・564ページ}
\kgbrow{当たる範囲}{中1〜中3の関数（比例と反比例・一次関数・$y=ax^2$）、データの活用（ヒストグラム・箱ひげ図・確率・標本調査）、数と式の一部（連立方程式・展開と因数分解・二次方程式）。数学Ⅰ（二次関数・二次不等式・データの分析）、数学Ⅱ（3次の展開・恒等式・判別式・直線の式・指数関数・対数）、数学Ⅲ（合成関数・逆関数）、数学A（組合せ・期待値）、数学B（正規分布・標本の平均のずれ）}
\kgbrow{目安とのずれ}{高校の内容が多い（3-3・3-5・3-8・3-10〜3-12・3-14〜3-17）。中学校の図形は第4巻にある}
\kgbrow{指導要領の外}{ブロック線図（3-3）、内挿・外挿（3-4）、片対数グラフ・デシベル（3-12）、最小二乗法・回帰直線（3-16。数学Bの「散布図に表したデータを関数とみなして処理する」は入口だけ）}
\end{kgbcard}

\begin{kgbcard}{4}{形と向き（図形・三角関数・ベクトル）}
\kgbstrip{4}{3,4,5,6,7,8,9,10,11,12,15,16}{8}{13}
\kgbrow{学年の目安}{中学2年〜高校2年　{\sffamily ★の範囲} \kgbst{1}{5}　16章・560ページ}
\kgbrow{当たる範囲}{小3〜小6の図形（面積・体積・合同・対称・拡大図と縮図・角柱と円柱）、中1〜中3の図形（作図をのぞく）、数学Ⅰの図形と計量（三角比・正弦定理・余弦定理）、数学Ⅱの三角関数と図形と方程式の一部（距離・円の式・内分点）、数学Aの図形の性質の一部、数学Cのベクトルと、平面上の曲線の一部（極座標・媒介変数・だ円）}
\kgbrow{目安とのずれ}{小3〜中1の図形（4-1・4-2・4-3・4-5、★〜★★）と、大学の内容（4-14・4-15、★★★★★）がある}
\kgbrow{指導要領の外}{外積（4-14）、円柱座標・球座標（4-15）、sin$^{-1}$ などで角を求めること（4-6）、ロボットの腕の順運動学・逆運動学（4-0・4-7）}
\end{kgbcard}

\begin{kgbcard}{5}{変化をつかむ（微分と積分）}
\kgbstrip{5}{11,12,13,14,16}{10}{13}
\kgbrow{学年の目安}{高校1年〜2年　{\sffamily ★の範囲} \kgbst{2}{4}　19章・492ページ}
\kgbrow{当たる範囲}{数学Ⅱの微分・積分の考え、数学Ⅲの極限・微分法・積分法（体積・曲線の長さは一部）、数学Bの $\Sigma$}
\kgbrow{目安とのずれ}{数学Ⅲは、一般的な例では高校3年。目安の「高校1年〜2年」より早い}
\kgbrow{指導要領の外}{テイラー展開（5-17）、1次近似と誤差の伝わり（5-8）、微分方程式の入口（5-13・5-18 の初期条件・一般解・特殊解）}
\end{kgbcard}

\begin{kgbcard}{6}{未来を計算する（微分方程式と複素数）}
\kgbstrip{6}{12,13,14,15,16}{12}{15}
\kgbrow{学年の目安}{高校2年〜3年　{\sffamily ★の範囲} \kgbst{2}{5}　16章・600ページ}
\kgbrow{当たる範囲}{数学Ⅱの複素数（6-8）、数学Cの複素数平面とド・モアブルの定理（6-9）}
\kgbrow{目安とのずれ}{序章を除く15章のうち13章が指導要領の外（大学の内容）。それでも★★の章（6-1・6-2と序章6-0）がある}
\kgbrow{指導要領の外}{微分方程式（6-0〜6-7、6-10〜6-15）、オイラーの公式（6-10）、双曲線関数（6-11）、数値解法（6-2）、指数分布とポアソン分布（6-4）}
\end{kgbcard}

\begin{kgbcard}{7}{まとめて計算する（行列）}
\kgbstrip{7}{13,14,15,16}{14}{16}
\kgbrow{学年の目安}{高校3年（発展）〜大学1年　{\sffamily ★の範囲} \kgbst{3}{5}　12章・492ページ}
\kgbrow{当たる範囲}{数学Cの (3)「離散グラフや行列を用いて工夫して表現することの意義」に当たる所から入る。7-10 に数学Bの等比数列・漸化式、数学Ⅲの無限等比級数}
\kgbrow{目安とのずれ}{指導要領では行列は「表現の意義」だけで、計算はすべて指導要領の外}
\kgbrow{指導要領の外}{行列の演算、一次変換、行列式、逆行列と掃き出し法、回転行列、固有値と固有ベクトル、対角化、使いみち（7-11：逆運動学・四端子定数・分布定数回路など）}
\end{kgbcard}

\begin{kgbcard}{8}{空間に広がる量（ベクトル解析）}
\kgbstrip{8}{14,15,16}{16}{16}
\kgbrow{学年の目安}{大学1年〜2年　{\sffamily ★の範囲} \kgbst{3}{5}　15章・632ページ}
\kgbrow{当たる範囲}{ほぼすべてが指導要領の外。わずかに、だ円の式（8-1）、媒介変数と曲線の長さ（8-7）が数学Ⅲ・数学Cに当たる}
\kgbrow{指導要領の外}{2変数の関数、偏微分、全微分、重積分、勾配、ベクトル場、線積分、面積分、発散、回転、ガウスの定理・ストークスの定理、ラプラシアン、熱方程式・波動方程式、マクスウェル方程式を読む（8-14 読みもの）}
\end{kgbcard}

\begin{kgbcard}{9}{波と変化を読みとく（フーリエ解析とラプラス変換）}
\kgbstrip{9}{12,14,16}{16}{16}
\kgbrow{学年の目安}{大学2年　{\sffamily ★の範囲} \kgbst{3}{5}　14章・640ページ}
\kgbrow{当たる範囲}{ほぼすべてが指導要領の外。部分分数の前の多項式の計算（9-7）、3次の特性方程式（9-9）が数学Ⅱの一部に、無限級数としての $e^{At}$（9-11）が数学Ⅲの一部に当たる}
\kgbrow{指導要領の外}{フーリエ級数・フーリエ変換・FFT、ラプラス変換・伝達関数・極と安定性、状態方程式、行列の指数関数、使いみち（9-12 読みもの）、全9巻の地図（9-13 終章）}
\end{kgbcard}

%% ====================================================================
\kgbsec{6}{追加した内容と、その理由}\label{kgb-ch6}

このシリーズには、現行の指導要領にない内容や、指導要領よりずっと詳しく扱う内容があります。この章では、それを4つに分けて示し、入れた理由を添えます。理由は、企画の依頼文、編集方針、委員会の記録、各巻の「物理・工学ではこう使う」に書かれていることだけを書きました。記録に個別の理由がないものは、その章の題材（どの工学で使うか）を書いています。

\par\Needspace{80mm}% 見出しだけがページの下に残らないように
\kgbsub{6.1}{大学の内容（第6〜9巻と、第4巻の一部）}

\par\addvspace{4pt}
\kgbcap{追加した内容（1）大学の内容}{}
\begin{bktable}[long, theme=kgb]{hline{Z}={0pt,white}, vline{1-Z}={Z}{0pt,white}, row{Z}={abovesep=3pt, belowsep=1pt}, hspan=minimal, colspec={Q[l,m,wd=30mm] Q[l,m,wd=21mm] Q[l,m,wd=23mm] Q[l,m,wd=52mm]}, rowhead=1, rowsep=1.5pt, colsep=3pt, cells={font=\footnotesize}, row{1}={font=\sffamily\bfseries\footnotesize}}
内容 & 章 & 指導要領 & 入れた理由 \\
微分方程式（立式・オイラー法・変数分離・振動・特性方程式・共振・連立） & 第6巻、5-13・5-18（入口）、9-5〜9-11 & 数学Ⅲの本文に記載なし & 企画の依頼文に「古典制御や現代制御におけるラプラス変換を用いた物理現象の微分方程式の立式や解析」。編集方針4「微分方程式は『次の瞬間を決めるルール』。式で解く前に表計算で解く」 \\
行列（演算・一次変換・行列式・逆行列・回転行列・固有値・対角化） & 第7巻、9-10・9-11 & 数学C (3) は「表現の意義」だけ & 依頼文に「動力学や逆運動学などの行列や行列式による計算、回路網理論における分布定数回路の四端子定数モデル」。委員会（機械工学）「ロボットには3次元のベクトル、外積、回転行列、逆運動学が要ります」 \\
オイラーの公式・フェーザ & 6-10、9-3・9-5 & 数学Cは極形式・ド・モアブルの定理まで & 路線④（角→sin→複素数）の終点「音・交流」へ向かう道具。題材：家のコンセントと交流の計算、三相交流の発電機 \\
双曲線関数 & 6-11 & なし & 題材：送電線のたるみ、アーチ（さかさまの懸垂線） \\
偏微分・全微分・重積分 & 8-1〜8-3・8-6 & なし & 依頼文に「電磁気学や流体力学などの3次元のベクトル解析」。委員会（数学）「専門書を読むのに欠かせない数学を足してください。全微分、重積分、ラプラシアン」 \\
ベクトル解析（勾配・発散・回転・線積分・面積分・ガウスとストークスの定理） & 8-4〜8-12、8-14 & 数学Cのベクトルは内積まで & 依頼文（同上）。委員会（電気工学）「式を見て『わき出しとうずの式だ』と言えるようにはしたい。専門書で一番つまずくのは、記号の読み方です」 \\
外積、円柱座標・球座標 & 4-14・4-15、8-6 & 「外積」の語はない & 第1回の委員会で「足すもの」に外積といろいろな座標を挙げた。題材：電気モーター、静止衛星 \\
偏微分方程式（熱方程式・波動方程式・ラプラス方程式） & 8-12・8-13、9-1 & なし & 路線①（変化）の終点「流体・熱」。題材：スマートフォンの熱をにがす、地震の波、楽器の弦 \\
フーリエ級数・フーリエ変換（FFT・サンプリング定理） & 9-1〜9-4、9-12 & 数学Ⅲは無限等比級数まで & 依頼文に「フーリエ級数展開など機械工学、電気工学などで必須の知識」 \\
ラプラス変換・伝達関数・極と安定性 & 9-5〜9-9 & なし & 依頼文（微分方程式の行と同じ）。委員会（電気工学）「ラプラス変換は『伝達関数』と『極と安定性』までにします」 \\
状態方程式・行列の指数関数 & 9-10・9-11 & なし & 委員会（電気工学）「状態方程式は、行列と微分方程式の話として残します」 \\
数値解法（オイラー法・差分法・掃き出し法・べき乗法・DFT と FFT） & 6-2・7-5・7-8・8-13・9-4 & 数値解法の項目はない（コンピュータの活用は各所にある） & 編集方針4・5。委員会（機械工学）「行列のかけ算もフーリエ係数も、表計算なら中学生が自分で出せます」 \\*
\SetCell[c=4]{c} \kgbcapline & & & \\
\end{bktable}
\par\addvspace{4pt}

\kgbsub{6.2}{行列の扱いの移り変わり}

行列は、指導要領の改訂のたびに扱いが変わってきた内容です。現行の指導要領では、行列は数学Cの (3)「数学的な表現の工夫」の中で「表現の道具」として挙げられているだけで、行列の和・積、逆行列、行列式、一次変換、固有値の語は本文にありません。工学では、ロボットの動き、回路網、制御、振動の解析など、行列の計算を使わない分野はほとんどありません。このシリーズが第7巻を丸ごと行列に当てたのは、このためです。

\begin{bkfigure}
\begin{tikzpicture}[x=1mm, y=1mm]
  \draw[bkGray, line width=1.2pt, -{Stealth[length=3mm]}] (2,0) -- (138,0);
  \foreach \kgbx/\kgbt/\kgbb in {
    22/平成11年告示/行列は{\sffamily\bfseries 数学C}の内容,
    70/平成21年告示/数学Cをなくす。行列は「{\sffamily\bfseries 数学活用}」で「図，表，行列及び離散グラフなどを用いて，事象を数学的に表現し考察すること」,
    118/平成30年告示（現行）/数学Cを新しく置く。(3)「数学的な表現の工夫」で「離散グラフや{\sffamily\bfseries 行列}を用いて工夫して表現することの意義」} {
    \fill[bkVol0] (\kgbx,0) circle (1.8);
    \node[kgb lab, anchor=south, font=\sffamily\bfseries\footnotesize] at (\kgbx,3) {\kgbt};
    \node[kgb gbox, anchor=north, text width=40mm] at (\kgbx,-4) {\kgbb};}
\end{tikzpicture}
\bkfigcaption{高校の指導要領での行列の扱い（文部科学省の資料による）}
\end{bkfigure}

\kgbnote{※ 文部科学省は、学校設定科目などで使える高等学校数学科教材「行列入門」を公開しています。現行の指導要領について、この資料は行列の扱いが「従前と比べて大きな変更はありません」としています。}

\kgbsub{6.3}{高校の範囲を超える内容（主に第3〜6巻）}

\par\addvspace{4pt}
\kgbcap{追加した内容（2）高校の範囲を超える内容}{}
\begin{bktable}[long, theme=kgb]{hline{Z}={0pt,white}, vline{1-Z}={Z}{0pt,white}, row{Z}={abovesep=3pt, belowsep=1pt}, hspan=minimal, rowhead=1, colspec={Q[l,m,wd=30mm] Q[l,m,wd=21mm] Q[l,m,wd=23mm] Q[l,m,wd=52mm]}, rowsep=1.5pt, colsep=3pt, cells={font=\footnotesize}, row{1}={font=\sffamily\bfseries\footnotesize}}
内容 & 章 & 指導要領 & 入れた理由 \\
テイラー展開、$n$ 回微分 & 5-17 & 数学Ⅲは無限等比級数まで。第$n$次導関数は用語にない & 第1回の委員会で「足すもの」に挙げた。題材：振り子とばね、トランジスタと「小信号」 \\
最小二乗法・回帰直線 & 3-16、7-5 & 数学Bは入口だけ & 委員会（数学）「統計を入れるなら、平均だけでは足りない。標準偏差、正規分布、最小二乗法まで」。（電気工学）「計測とノイズを扱うには、確率と統計が欠かせません」 \\
指数分布・ポアソン分布 & 6-4 & なし & 題材：放射性炭素の年代 \\
1次近似・誤差の伝わり & 5-8・8-3 & 項目なし & 題材：カーナビの推測航法。5-8 は 6-2 のオイラー法の予告 \\
片対数グラフ・デシベル & 3-12 & なし & 題材：ダイオードの電流、ボード線図 \\
ブロック線図 & 3-3・9-8 & なし & 題材：温度センサー（温度を入れると電圧が出てくる箱） \\
sin$^{-1}$ などで角を求める & 4-6 & なし & 題材：すべり出す角で摩擦をはかる \\
順運動学・逆運動学 & 4-0・4-7・7-6・7-11 & なし（工学） & 委員会（機械工学）「余弦定理は逆運動学の準備として入れてください」 \\*
\SetCell[c=4]{c} \kgbcapline & & & \\
\end{bktable}
\par\addvspace{4pt}

\kgbsub{6.4}{算数・数学の外にある土台}

\par\addvspace{4pt}
\kgbcap{追加した内容（3）算数・数学の外にある土台}{}
\begin{bktable}[long, theme=kgb]{hline{Z}={0pt,white}, vline{1-Z}={Z}{0pt,white}, row{Z}={abovesep=3pt, belowsep=1pt}, hspan=minimal, rowhead=1, colspec={Q[l,m,wd=30mm] Q[l,m,wd=21mm] Q[l,m,wd=23mm] Q[l,m,wd=52mm]}, rowsep=1.5pt, colsep=3pt, cells={font=\footnotesize}, row{1}={font=\sffamily\bfseries\footnotesize}}
内容 & 章 & 指導要領 & 入れた理由 \\
量＝数値×単位、SI単位・組立単位・接頭語、単位つきの換算 & 2-1・2-3・2-10 & 小3で k・m に「触れる」程度 & 編集方針1「量＝数値×単位から出発する」。題材：単位のとりちがえで消えた火星探査機 \\
有効数字・誤差（±）・ずれの割合 & 2-6、3-14 & 中3の「誤差、近似値」 & 路線⑦（目もり→誤差→統計）。題材：部品の寸法と公差 \\
変化量 $\Delta$・1あたりの量・量の掛け算 & 2-3 & $\Delta$ はどの学年の用語・記号にもない & 編集方針3「大学まで続く種を早くまく（ひき算＝変化量、かけ算＝面積…）」 \\
表計算・プログラム & 2-18・2-19（以後の全巻で使う） & 「コンピュータなどを用いて」（各所）、小5のプログラミングの例 & 編集方針5。第1回の委員会で「表計算とプログラム」を足すと決めた \\
16進数（2進数は 6.5） & 2-17 & 16進数はない & 題材：センサーの値を数にする（分解能）、マイコンボード \\*
\SetCell[c=4]{c} \kgbcapline & & & \\
\end{bktable}
\par\addvspace{4pt}

\kgbsub{6.5}{指導要領と同じ項目を、詳しく扱う所}

\par\addvspace{4pt}
\kgbcap{追加した内容（4）同じ項目を、詳しく扱う所}{}
\begin{bktable}[long, theme=kgb]{hline{Z}={0pt,white}, vline{1-Z}={Z}{0pt,white}, row{Z}={abovesep=3pt, belowsep=1pt}, hspan=minimal, rowhead=1, colspec={Q[l,m,wd=24mm] Q[l,m,wd=14mm] Q[l,m,wd=42mm] Q[l,m,wd=46mm]}, rowsep=1.5pt, colsep=3pt, cells={font=\footnotesize}, row{1}={font=\sffamily\bfseries\footnotesize}}
項目 & 章 & 指導要領の扱い & このシリーズの扱い \\
反比例 & 2-15・3-7 & 小6は「知る」程度 & $xy$ が一定、グラフ、$pV$ が一定 \\
三平方の定理 & 4-4 & 中3は「証明できることを知る」 & 2通りの証明を本文で見せる \\
正弦定理・余弦定理 & 4-7 & 数学Ⅰ & ロボットの腕の順運動学・逆運動学、力の三角形 \\
加法定理 & 4-9 & 数学Ⅱ & 2通りで導く、積と和の式、うなり、sin$^2$ の平均（実効値 5-15） \\
内積 & 4-12 & 数学C & 仕事、正射影、相関係数は偏差ベクトルの向きのそろいぐあい \\
微分 & 5-4 & 数学Ⅱは三次までの関数が中心、極限は直観的に & $x^n$ の一般の形、表から規則を見つける \\
積分 & 5-10 & 数学Ⅱは二次までの関数が中心 & 区分求積・$\Sigma$ から定積分へ \\
標準偏差・相関 & 3-14・3-15 & 数学Ⅰ & 共分散、$\sigma/\sqrt{n}$、単位を変えたときの $\sigma$ \\
二次方程式 & 3-10 & 中3は実数解をもつものだけ & 判別式、2次不等式。虚数の解は 6-8 \\
指数・対数 & 3-11・3-12 & 数学Ⅱ & 片対数グラフ、デシベル、半減期 \\
2進数 & 2-17 & 数学Aの内容の取扱いで、数学と人間の活動の題材の例に「二進法」 & 2進数と10進数の行き来、ビットの数と何通り（$2^n$）、バイト、分解能。16進数も扱う（16進数は指導要領の外：6.4） \\*
\SetCell[c=4]{c} \kgbcapline & & & \\
\end{bktable}
\par\addvspace{4pt}

\par\Needspace{165mm}% 見出しから図までを同じページに
\kgbsub{6.6}{まとめ：路線ごとに見た、指導要領の先}

6.1〜6.5 の内容を、7本の路線ごとにまとめます。指導要領の中でどこまで進むかと、このシリーズでどこまで進むかを並べると、次のようになります。

\begin{bkfigure}
\begin{tikzpicture}[x=1mm, y=1mm]
  \node[kgb lab, font=\sffamily\bfseries\footnotesize] at (49,6) {指導要領の中};
  \node[kgb lab, font=\sffamily\bfseries\footnotesize] at (109,6) {このシリーズで、その先へ};
  \foreach \kgbr/\kgbn/\kgba/\kgbb/\kgbe [count=\kgbi from 0] in {
    1/変化/ひき算→変化の割合・傾き（中2）\\→微分（数Ⅱ・数Ⅲ）/微分方程式（第6巻）\\→偏微分・勾配・熱方程式（第8巻）/流体・熱,
    2/たし合わせ/かけ算→面積（小4〜）\\→積分（数Ⅱ・数Ⅲ）/重積分・線積分・面積分、\\ガウスの定理（第8巻）/電磁気,
    3/倍々/累乗→指数関数・対数関数（数Ⅱ）\\→$e$（数Ⅲ）/時定数（6-4）\\→ラプラス変換・伝達関数（第9巻）/制御,
    4/回転と波/角（小4）→三角関数（数Ⅱ）\\→複素数平面（数C）/オイラーの公式（6-10）\\→フーリエ解析（第9巻）/音・交流,
    5/表/表とグラフ（小2〜）\\→行列で表現する意義（数C）/行列の計算・固有値（第7巻）\\→状態方程式（第9巻）/回路網・制御,
    6/矢印/数直線・座標\\→ベクトルと内積（数C）/外積（4-14）\\→ベクトル解析（第8巻）/ロボット・電磁気,
    7/はかる/目もり→データの分析（数Ⅰ）\\→統計的な推測（数B）/最小二乗法（3-16）\\→スペクトル・FFT（9-4）/計測・診断} {
    \pgfmathsetmacro\kgby{-\kgbi*18.5}
    \node[anchor=west, inner sep=0pt] at (0,\kgby) {\bkrouteicon[1.3em]{\kgbr}};
    \node[kgb lab, anchor=west, font=\sffamily\bfseries\footnotesize] at (5.8,\kgby) {\kgbn};
    \node[kgb gbox, anchor=west, text width=44mm, minimum height=15.5mm] (g\kgbi) at (25,\kgby) {\kgba};
    \node[kgb tbox, anchor=west, text width=46mm, minimum height=15.5mm] (t\kgbi) at (82,\kgby) {\kgbb};
    \draw[kgb arr] (g\kgbi) -- (t\kgbi);
  }
\end{tikzpicture}
\bkfigcaption{路線ごとの、指導要領の中の到達点と、このシリーズの到達点}
\end{bkfigure}


%% ====================================================================
\kgbsec[含めていない・軽く扱う内容と、補い方]{7}{含めていない・軽く扱う内容と、\\補い方}\label{kgb-ch7}

このシリーズは、物理・工学で使う道筋を優先したため、指導要領にあっても含めていない内容、軽く扱っている内容があります。正直にお示しし、学校の教科書・問題集での補い方を添えます。

\kgbsub{7.1}{方針として軽くした内容}

次の内容は、編集委員会が構成を決めるときに、方針として軽くしたものです。「理由」は委員会の記録の言葉です。

\par\addvspace{4pt}
\kgbcap{方針として軽くした内容}{}
\begin{bktable}[long, theme=kgb]{hline{Z}={0pt,white}, vline{1-Z}={Z}{0pt,white}, row{Z}={abovesep=3pt, belowsep=1pt}, hspan=minimal, rowhead=1, colspec={Q[l,m,wd=27mm] Q[l,m,wd=25mm] Q[l,m,wd=42mm] Q[l,m,wd=32mm]}, rowsep=1.5pt, colsep=3pt, cells={font=\footnotesize}, row{1}={font=\sffamily\bfseries\footnotesize}}
内容 & 指導要領 & このシリーズでの扱い & 理由 \\
図形の証明・作図 & 小3（作図）、中1（作図）、中2・中3（証明）、数学A（作図） & 証明は 4-1 の三角形の内角の和、4-4 の三平方の定理（2通り）、4-9 の加法定理の導き方。作図は 4-2 で垂直二等分線を紙を折って見つける程度。作図の手順の節はない & 定理の意味が分かれば十分 \\
整数の性質（倍数判定・不定方程式・互除法） & 小5、中1、数学A & 約分に必要な分だけ（2-4：素数・公約数・公倍数）。偶数・奇数の定義は 2-9 の「もっと知りたい」 & 物理・工学でほとんど使わない \\
場合の数（順列・組合せ） & 数学A & 確率に必要な範囲だけ（3-17：樹形図・パスカルの三角形・組合せ） & 計算演習より、ばらつきの見方が大事 \\
積分の計算演習（回転体の体積・曲線の長さ） & 数学Ⅲ & 積分の意味（面積・仕事・流量）にしぼる。体積は 5-10 の例題、曲線の長さは 6-11・8-7 に式がある & 計算は数式処理ソフトで確かめられる \\
集合と論理 & 数学Ⅰ & 4-1 の「逆と反例」（仮定・結論）だけ。集合は扱わない & 必要最小限 \\
数列の一般論（等差数列・数学的帰納法） & 数学B & $\Sigma$ は 5-10、漸化式は 6-2（オイラー法の形）と 7-10。等差数列・帰納法はない & 微分方程式の数値解に吸収できる \\*
\SetCell[c=4]{c} \kgbcapline & & & \\
\end{bktable}
\par\addvspace{4pt}

\kgbsub{7.2}{今の版で扱っていない、または一部だけの内容}

このシリーズは、物理や工学で使わない単元は省くか軽くする方針（編集方針7）で作っており、次の内容は扱っていません（一部だけ扱うものを含みます）。補い方の欄の章を読んだ後に、学校の教科書・問題集で学んでください（図\ref{kgb-f-hokai}）。特に、掛け算・割り算の筆算は、今の版の第1巻では扱っていません。第1巻を使うときは、小学校の教科書で筆算の練習をしてください。

\par\addvspace{4pt}
\kgbcap{今の版で扱っていない、または一部だけの内容}{}
\begin{bktable}[long, theme=kgb]{hline{Z}={0pt,white}, vline{1-Z}={Z}{0pt,white}, row{Z}={abovesep=3pt, belowsep=1pt}, hspan=minimal, colspec={Q[l,m,wd=33mm] Q[l,m,wd=20mm] Q[l,m,wd=56mm] Q[l,m,wd=17mm]}, rowhead=1, rowsep=1.5pt, colsep=3pt, cells={font=\footnotesize}, row{1}={font=\sffamily\bfseries\footnotesize}}
内容 & 指導要領 & このシリーズでの扱い & 補い方 \\
掛け算・割り算の筆算 & 小3・小4 & 筆算は、足し算・引き算（1-7）だけ。掛け算は九九（1-12・1-13）と「3×10」程度、割り算は九九の範囲（1-14・1-15）。23×14 の筆算は 3-8 のコラムで「展開」として見せるだけ & 1-13・1-15 の後 \\
そろばん & 小3・小4 & 扱っていない & 1-15 の後 \\
円グラフ・帯グラフ & 小5 & 扱っていない & 2-5 の後 \\
独立な試行・条件付き確率 & 数学A & 扱っていない & 3-17 の後 \\
統計的な推測（確率分布・二項分布の名前、区間推定、信頼区間、有意水準） & 数学B & 正規分布と標本調査（3-17）、$\sigma/\sqrt{n}$（3-14）。仮説検定は 3-17 の「もっと知りたい」だけ。区間推定はない & 3-17 の後 \\
三次の因数分解・因数定理・高次方程式 & 数学Ⅱ & 3-8 で $(a+b)^3$ の展開、9-9 で3次の特性方程式を扱う程度 & 3-8 の後 \\
多項式の除法・分数式 & 数学Ⅱ & 9-7 で部分分数の前に1行使うだけ & 3-8 の後（9-7 の前に） \\
等式・不等式の証明 & 数学Ⅱ & 扱っていない（恒等式・係数比較は 3-8） & 3-8 の後 \\
二項定理 & 数学Ⅱ（用語） & 語はない。パスカルの三角形（3-8・5-4 の「もっと知りたい」、3-17） & 3-8 の後 \\
外分点・軌跡・不等式の表す領域 & 数学Ⅱ & 扱っていない（内分点は 4-11） & 4-11 の後 \\
分数関数・無理関数のグラフ & 数学Ⅲ & $y=a/x$（3-7）だけ & 3-7 の後 \\
凹凸・変曲点 & 数学Ⅲ & 5-7「2回微分で、山と谷を見分ける」まで。変曲点の語はない & 5-7 の後 \\
三角形・円の性質（内心・メネラウスの定理・チェバの定理・方べき・接弦定理など） & 数学A & 外心（4-2）、重心（4-11）、内接円（4-7 の「もっと知りたい」）だけ & 4-2 の後 \\
二次曲線（焦点・準線） & 数学C & 放物線（3-9）、だ円（4-8 の「もっと知りたい」・4-15・8-1）、双曲線は $y=a/x$ のグラフとして（3-7）。焦点・準線で3つをまとめる扱いはない & 4-15 の後 \\
数学と人間の活動・数学と社会生活 & 数学A・数学B の (3) & 2-17（2進数）、3-16（最小二乗法）が近い内容 & 学校の授業で \\
逆三角関数の微分・不定形の極限（ロピタルの定理） & 大学（微分積分学） & 扱っていない（sin$^{-1}$ は 4-6 で電卓を使って角を求めるだけ） & 5-15 の後に大学の教科書で \\*
\SetCell[c=4]{c} \kgbcapline & & & \\
\end{bktable}
\par\addvspace{4pt}

\par\Needspace{160mm}% 見出しと図を同じページに
\kgbsub{7.3}{補い方：どの章の後に学ぶとよいか}

7.1・7.2 の内容は、学校の教科書・問題集で補うことができます。次の図は、このシリーズのどの章を読んだ後に補うとつながりやすいかの例です。かっこの中は、指導要領で学ぶ学年・科目です。

\begin{bkfigure}
\begin{tikzpicture}[x=1mm, y=1mm]
  \draw[bkGray, line width=1.6pt] (11,2) -- (11,-133.4);
  \foreach \kgby/\kgbc/\kgbv/\kgbt [count=\kgbi from 0] in {
    0.0/1-13/1/掛け算の筆算（小3）,
    -8.6/1-15/1/割り算の筆算（小4）、そろばん（小3・小4）,
    -17.2/2-4/2/最大公約数・最小公倍数・素因数分解の名前、偶数と奇数（小5・中1）,
    -25.8/2-5/2/円グラフ・帯グラフ（小5）,
    -34.4/3-7/3/分数関数・無理関数のグラフ（数学Ⅲ）,
    -45.8/3-8/3/三次の因数分解・多項式の除法・分数式・因数定理・二項定理・等式と不等式の証明（数学Ⅱ）,
    -59.8/3-17/3/順列・和の法則と積の法則・独立な試行・条件付き確率（数学A）、確率分布・区間推定（数学B）,
    -71.2/4-1/4/図形の証明・作図（中1〜中3・数学A）、集合と命題（数学Ⅰ）,
    -79.8/4-2/4/三角形・円の性質（内心・方べきなど）（数学A）,
    -88.4/4-11/4/外分点・軌跡・不等式の表す領域（数学Ⅱ）,
    -97.0/4-15/4/二次曲線の焦点・準線（数学C）,
    -105.6/5-7/5/凹凸・変曲点（数学Ⅲ）,
    -114.2/5-10/5/等差数列・数学的帰納法（数学B）,
    -122.8/5-15/5/逆三角関数の微分・不定形の極限（大学の微分積分学）,
    -131.4/5-16/5/回転体の体積・曲線の長さの計算練習（数学Ⅲ）} {
    \node[fill=bkVol\kgbv, text=bkOn\kgbv, rounded corners=1mm, font=\sffamily\bfseries\footnotesize, minimum width=12mm, inner ysep=1mm] (c\kgbi) at (11,\kgby) {\kgbc};
    \node[kgb gbox, anchor=west, text width=106mm, minimum height=7mm, inner ysep=1.1mm] (b\kgbi) at (27,\kgby) {\kgbt};
    \draw[kgb arr] (c\kgbi) -- (b\kgbi);}
  \node[kgb lab, anchor=south, font=\sffamily\bfseries\footnotesize] at (11,4) {このシリーズの章};
  \node[kgb lab, anchor=south west, font=\sffamily\bfseries\footnotesize] at (27,4) {その後に、学校の教科書・問題集で補う内容};
\end{tikzpicture}
\bkfigcaption{補い方の例（章を読んだ後に、学校の教科書・問題集で学ぶ）}\label{kgb-f-hokai}
\end{bkfigure}

多項式の割り算は、第9巻 9-7 でラプラス変換の部分分数に使います。9-7 を読む前に、学校の教科書で学んでおくと安心です。


\kgbsub{7.4}{入試との関係}

このシリーズは、入試を解くための問題演習を目的にしていません（1.1、企画の依頼文）。高校入試・大学入試の準備には、この章の内容を含め、学校の教科書・問題集と併せて使ってください。

なお、令和7年度の大学入学共通テストの数学では、『数学Ⅰ，数学Ａ』の数学Aは「図形の性質」「場合の数と確率」の2項目から、『数学Ⅱ，数学Ｂ，数学Ｃ』は数学Bの「数列」「統計的な推測」と数学Cの「ベクトル」「平面上の曲線と複素数平面」の4項目から3項目を選んで解答する形で出題され、数学A・B・Cの (3) の項目は出題されません。7.1・7.2 の内容のうち、図形の証明・作図、三角形・円の性質、場合の数、独立な試行・条件付き確率、数列、統計的な推測、二次曲線は、これらの出題範囲の項目に含まれます。

%% ====================================================================
\kgbsec{8}{大学の工学部の数学との対応}\label{kgb-ch8}

このシリーズの第3巻の統計の章、第4巻の後半、第5〜9巻は、大学の工学部の1〜2年で普通学ぶ数学の土台になっています。科目の名前・中身・学年は大学や学科ごとに違うので、ここでは実際のシラバスの例から拾った「よくある範囲」と照らしました（表\ref{kgb-t-univ}・図\ref{kgb-f-univ}）。


\par\addvspace{4pt}
\kgbcap{大学の工学部の数学の科目と、このシリーズの章}{kgb-t-univ}
\begin{bktable}[long, theme=kgb]{hline{Z}={0pt,white}, vline{1-Z}={Z}{0pt,white}, row{Z}={abovesep=3pt, belowsep=1pt}, hspan=minimal, colspec={Q[l,m,wd=24mm] Q[l,m,wd=38mm] Q[l,m,wd=30mm] Q[l,m,wd=34mm]}, rowhead=1, rowsep=1.5pt, colsep=3pt, cells={font=\footnotesize}, row{1}={font=\sffamily\bfseries\footnotesize}}
大学の科目 & 代表的な内容（シラバスの例） & このシリーズの章 & このシリーズで扱わないこと \\
微分積分学（1変数） & 極限、逆三角関数とその微分、高次導関数、テイラーの定理、積分法 & 第5巻（5-3〜5-18）、テイラー展開 5-17 & $\varepsilon$-$\delta$ 論法は使わない（極限は、数値で近づくようすで示す）。逆三角関数の微分、不定形の極限（ロピタルの定理）は扱わない（7.2） \\
微分積分学（多変数） & 偏微分、全微分、ヤコビアン、極値問題、重積分 & 8-1〜8-4・8-6、ヤコビ行列 7-2 & ― \\
線形代数学Ⅰ & 行列の演算、行列式、連立一次方程式（掃き出し法）、逆行列 & 7-1・7-3〜7-5 & ― \\
線形代数学Ⅱ & ベクトル空間、線形写像、固有値・固有ベクトル、対角化、直交行列 & 7-2・7-6〜7-10、9-10・9-11 & 抽象的なベクトル空間（行列は2×2と3×3の具体的なもの） \\
常微分方程式 & 変数分離形、一階線形、高階線形・連立線形 & 第6巻、5-13・5-18（入口）、9-7・9-10・9-11 & ― \\
ベクトル解析 & 勾配・発散・回転、線積分・面積分、ガウス・ストークスの定理、直交曲線座標 & 4-14・4-15、8-4〜8-12・8-14 & ― \\
偏微分方程式 & 波動方程式・熱伝導方程式・ラプラス方程式（変数分離・フーリエ級数による解法） & 8-12・8-13、9-1 & ― \\
複素関数論 & 正則関数、コーシーの積分定理、留数定理 & 6-8〜6-10（複素数とオイラーの公式まで） & 複素関数論（留数定理など）。ラプラス逆変換は変換表と部分分数で行う \\
フーリエ解析 & フーリエ級数、複素フーリエ級数、フーリエ変換 & 9-1〜9-4・9-12 & ― \\
ラプラス変換 & ラプラス変換の性質、微分方程式への応用 & 9-5〜9-9 & ― \\
確率統計 & 確率変数、二項・ポアソン・正規分布、推定と検定 & 3-13〜3-17、6-4 & 区間推定（仮説検定は 3-17 の「もっと知りたい」だけ） \\*
\SetCell[c=4]{c} \kgbcapline & & & \\
\end{bktable}
\par\addvspace{4pt}

\kgbnote{※ 「代表的な内容」は、東北大学・京都工芸繊維大学・北海道大学・神戸大学・東京農工大学のシラバスなどから拾った例です（出典の U1〜U6・U8・U9）。確率統計は、工学部向けのシラバスの例を確かめられなかったため、数学科向けのシラバスを目安にしました。}
\kgbnote{※ 第9巻の終章（9-13）にも、電気回路・制御工学・信号処理・機械力学・流体力学・電磁気学などの専門の科目と、その土台になる巻・章の表があります。}

\begin{bkfigure}
\begin{tikzpicture}[x=1mm, y=1mm]
  \foreach \kgbv [count=\kgbk] in {3,...,9} {\node[inner sep=0pt] at ({57+(\kgbk-1)*12.6},4) {\kgbvol{\kgbv}};}
  \foreach \kgbs/\kgbl/\kgbf [count=\kgbi from 0] in {
    微分積分学（1変数）/1年前期/{5/5},
    微分積分学（多変数）/1年後期/{7/4,8/5},
    線形代数学Ⅰ/1年前期/{7/5},
    線形代数学Ⅱ/1年後期/{7/5,9/4},
    常微分方程式/1〜2年/{5/4,6/5,9/5},
    ベクトル解析/2年/{4/4,8/5},
    偏微分方程式/2年/{8/5,9/4},
    複素関数論/2年/{6/4},
    フーリエ解析/2年/{9/5},
    ラプラス変換/2年/{9/5},
    確率統計/1〜2年/{3/5,6/4}} {
    \pgfmathsetmacro\kgby{-2-\kgbi*7.6}
    \ifodd\kgbi\else\fill[bkPaper] (0,{\kgby-3.8}) rectangle (140,{\kgby+3.8});\fi
    \node[kgb lab, anchor=west, font=\sffamily\bfseries\footnotesize] at (1,\kgby) {\kgbs};
    \node[kgb lab, anchor=east, text=bkGray] at (51,\kgby) {\kgbl};
    \foreach \kgbv/\kgbm in \kgbf {
      \pgfmathsetmacro\kgbx{57+(\kgbv-3)*12.6}
      \ifnum\kgbm=5 \fill[bkVol\kgbv, draw=bkInk, line width=0.6pt] (\kgbx,\kgby) circle (2.4);
      \else \draw[bkVol\kgbv!80!black, line width=1.4pt, fill=white] (\kgbx,\kgby) circle (2.1);\fi}
  }
  \fill[bkVol0, draw=bkInk, line width=0.6pt] (60,-89) circle (2.4);
  \node[kgb lab, anchor=west] at (64,-89) {土台の中心};
  \draw[bkVol0, line width=1.4pt, fill=white] (92,-89) circle (2.1);
  \node[kgb lab, anchor=west] at (96,-89) {一部・入口};
\end{tikzpicture}
\bkfigcaption{大学の工学部の数学の科目と、その土台になる巻（学年はよくある例）}\label{kgb-f-univ}
\end{bkfigure}

%% ====================================================================
\kgbsec{9}{★と学年の目安の読み方}\label{kgb-ch9}

\kgbsub{9.1}{★は「このシリーズの中での難しさ」の目安}

各章と「できるようになること」の各項目には、★を1〜5個付けています。★は、このシリーズの中での難しさ、つまり、その章でどの道具を使うかの目安です。編集を始めるとき、★の数を次のように決めました。これは、おおよその目安です。

\begin{bkcenter}
\begin{bktable}{colspec={Q[l,m,wd=30mm] Q[l,m,wd=46mm] Q[l,m,wd=40mm]}, rowsep=1.5pt, colsep=4pt, cells={font=\small}, row{1}={font=\sffamily\bfseries\small}}
★ & 内容の高さ（おおよその目安） & 子どもへの言い方（山登り） \\
\bkstars{1} & 小4まで & はじめの一歩 \\
\bkstars{2} & 小5〜中1 & ふもとを歩く \\
\bkstars{3} & 中2〜高1 & 山を登る \\
\bkstars{4} & 高2〜高3 & 尾根を歩く \\
\bkstars{5} & 大学（工学部） & 頂上をめざす \\
\end{bktable}
\bktabcaption{★の数と、内容の高さのおおよその目安}
\end{bkcenter}

ただし、★は指導要領の学年と一対一には対応していません。同じ指導要領の内容でも、このシリーズの書き方でたどれる高さに合わせて★を付けたため、目安の表より低い★が付いている章があります。たとえば、数学Ⅱ（一般的な例では高校2年）の内容は、目安の表どおりなら★★★★ですが、次の図のように★★〜★★★の章がいくつもあります。反対に、目安の表より高い★が付いている章は見つかりませんでした。なお、章の★が低くても、「できるようになること」の項目の★は目安の表どおりの章があります（4-2 の弧度法の項目は★★★★。9.4）。

\begin{bkfigure}
\begin{tikzpicture}[x=1mm, y=1mm]
  \foreach \kgbs in {1,...,5} {\draw[bkLine, line width=0.4pt] (18,{\kgbs*9}) -- (138,{\kgbs*9});
    \node[anchor=east, inner sep=0pt, font=\sffamily\footnotesize] at (17,{\kgbs*9}) {\bkstars{\kgbs}};}
  \fill[bkVol0!10] (18,{4*9-3.5}) rectangle (138,{4*9+3.5});
  \node[kgb lab, anchor=south east, text=bkVol0, font=\sffamily\bfseries\footnotesize] at (137.5,{4*9+3.7}) {目安の表で、数学Ⅱ（高2の例）に当たる★};
  \foreach \kgbc/\kgbv/\kgba/\kgbb/\kgbt [count=\kgbi from 0] in {
    2-7/2/1/2/{0乗・負\\の指数},3-11/3/2/3/{指数\\関数},3-12/3/2/3/対数,4-8/4/3/4/{三角\\関数},4-9/4/4/4/{加法\\定理},
    5-1/5/2/2/{平均の\\速さ},5-2/5/2/2/{極限の\\入口},5-3/5/3/3/{微分\\係数},5-9/5/2/2/{面積と\\道のり},6-8/6/4/4/複素数} {
    \pgfmathsetmacro\kgbx{26+\kgbi*10.8}
    \draw[bkVol\kgbv!85!black, line width=4pt, line cap=round] (\kgbx,{\kgba*9}) -- (\kgbx,{\kgbb*9+0.01});
    \fill[bkVol\kgbv, draw=bkInk, line width=0.5pt] (\kgbx,{\kgba*9}) circle (1.9) (\kgbx,{\kgbb*9}) circle (1.9);
    \node[fill=bkVol\kgbv, text=bkOn\kgbv, rounded corners=0.8mm, font=\sffamily\bfseries\footnotesize, inner xsep=0.8mm, inner ysep=0.5mm] at (\kgbx,0) {\kgbc};
}
\end{tikzpicture}
\bkfigcaption{数学Ⅱの内容を中心に扱う章と、その章の★（線の長さは★の範囲）}
\end{bkfigure}
\kgbnote{※ 図の章と、数学Ⅱに当たる内容：2-7 0乗・負の指数／3-11 指数関数／3-12 対数／4-8 三角関数／4-9 加法定理／5-1 平均の速さ（平均変化率）／5-2 極限の入口／5-3 微分係数・導関数／5-9 面積と道のり／6-8 複素数}

このように★が低くなっているのは、編集方針8にあるように、図・実験・表計算でたどれる書き方を選んでいるためです。★は「学校の何年生の内容か」ではなく、「このシリーズの中で、どれくらいの道具を使うか」と読んでください。学校の学年との対応は、4章の図と表、付録A・付録Bで確かめてください。

\kgbsub{9.2}{3つの段の章の★}

6-10〜6-14、7-7〜7-10、8-1〜8-13、9-1〜9-11 の各章（第6〜9巻の★★★★★の章は、読みものと終章を除いてすべて含まれます）は「目でわかる→手でわかる→式でわかる」の3段で書かれていて（1.5、図\ref{kgd-f-dan}）、「できるようになること」の3つの項目の★は、普通、段ごとに★★★・★★★★・★★★★★と並んでいます。章の帯に★★★★★とあっても、前の2段は★★★以下の道具（図・実験・表計算・プログラム）だけでたどれます。大学の内容の章を中学生・高校生が読むときは、「手でわかる」の終わりの「ここまでで十分」の旗まで読めれば、その章のゴールに着いたと考えてください。

\kgbsub{9.3}{学年の目安の読み方}

4章の「学年の目安」は、編集のときに決めた目安です。巻の中には、その学年より前や後に学校で習う内容も含まれています（図\ref{kgb-f-grid}）。学年の目安は「その学年になったら、この巻を読む」という決まりではありません。お子さんが学校でいま何を習っているかと、付録Aの対応表から、読み始める巻と章を選んでください。生徒用の紙面には学年を書いていないので、学年の目安より早く読んでも、遅く読んでも、子ども自身は気にせずに読めるようになっています。

\kgbsub{9.4}{照らし合わせるときの注意}

\kgbitem{{\sffamily\bfseries 1つの章が、いくつかの学年にまたがります。}★が範囲になっている章が34あります（例：2-2 ★〜★★、3-8 ★★〜★★★）。★が1つの章でも、1-2（小1〜小3）、2-6（小4〜中3）、3-13（小6〜数学Ⅰ）、4-2（小3〜数学C）のように、いくつかの学年にまたがる章が多くあります。章の中の「できるようになること」の各項目にも★が付いているので、項目ごとに照らすと正確です。}
\kgbitem{{\sffamily\bfseries 章の★より高い項目があります。}3-5（項目に★★★）、4-2（項目に★★★★：ラジアン）、5-8（項目に★★★）、6-2（項目に★★★）。}
\kgbitem{{\sffamily\bfseries 「もっと知りたい」は発展です。}章より★が高いことが多く、章の本体とは分けて数えています。}
\kgbitem{{\sffamily\bfseries 物理・工学の題材は、算数・数学の単元そのものではありません。}「物理・工学ではこう使う」のオームの法則（2-15）やコンデンサーの回路（6-4）などは理科の内容に当たります。理科の指導要領との対応は、この本では扱っていません。}

%% ====================================================================
\kgbsec{10}{授業・家庭での使い方の例}\label{kgb-ch10}

\kgbsub{10.1}{ご家庭で使うときに}

ノート・電卓・表計算・巻末の使い方は、第1章（1.8・1.9）にまとめました。ここでは、そばで見守る方へのお願いを挙げます。

\kgbitem{例題と「じぶんでとこう」で、章末問題に出てくる問題の形を、すべて1回ずつ見せています。章末問題で手が止まったら、同じ形の例題に戻るよう声をかけてください。}
\kgbitem{1回で分からなくても大丈夫です。章の「ふりかえり」の各項目に、戻るページが書いてあります。}
\kgbitem{学年の目安より早く読んでも遅く読んでも構いません（9.3）。生徒用の紙面には学年を書いていないので、「何年生の内容か」を本人に伝える必要はありません。}
\kgbitem{7章の内容（掛け算・割り算の筆算など）は、学校の教科書・問題集で補ってください（図\ref{kgb-f-hokai}）。}

\kgbsub{10.2}{学校で使うときに}

\kgbitem{章の流れはどの章も同じです（はてな？→やってみよう→数学の道具→例題→じぶんでとこう→まちがえやすいところ→物理・工学ではこう使う→章末問題）。授業の1時間を、この流れの一部に当てることができます。}
\kgbitem{「まちがえやすいところ」は、×と○を1行ずつ並べ、なぜ間違いかを書いています。授業でのつまずきの確認に使えます。}
\kgbitem{用語・書き方が学校の教科書と違う所（表\ref{kgb-t-yougo}）があります。特に虚数単位を $j$ と書く点（6-8 から）は、生徒に伝えてください。}

\par\Needspace{100mm}% 見出しと表10-1を同じページに
\kgbsub{10.3}{目的に合わせた使い方}

\par\addvspace{4pt}
\kgbcap{目的に合わせた使い方の例}{}
\begin{bktable}[long, theme=kgb]{hline{Z}={0pt,white}, vline{1-Z}={Z}{0pt,white}, row{Z}={abovesep=3pt, belowsep=1pt}, hspan=minimal, rowhead=1, colspec={Q[l,m,wd=32mm] Q[l,m,wd=94mm]}, rowsep=2pt, colsep=4pt, cells={font=\footnotesize}, row{1}={font=\sffamily\bfseries\footnotesize}}
目的 & 使い方の例 \\
学校の授業と並べて使う & 付録Aで、いま学校で習っている単元の章を探します。章の初めの「はてな？」と、終わりの「物理・工学ではこう使う」を先に読むと、その単元が何に使われるかが分かります \\
先へ進む & 巻の順に読むのが基本です。章の帯の下の「道しるべ」に、前に読む章とこの先の章が書いてあります。自分の読める★まで読んで、先へ進んでもかまいません \\
学び直す & 第1章の「どこから始めるか」（1.7）で、始める巻を選びます。分からない所は、道しるべの「この章の前に」の章へ戻ります \\
大学の工学に備える & 8章の表で、学ぶ科目の土台になる巻を選びます。3つの段の章は「ここまでで十分」の旗まで読み、その後「式でわかる」で専門書の書き方に慣れます \\
理科の授業とつなぐ & 「物理・工学ではこう使う」や「やってみよう」の実験は、理科の授業の導入や発展に使えます（例：2-15 のばねとオームの法則、6-4 のコンデンサーの放電） \\*
\SetCell[c=2]{c} \kgbcapline & \\
\end{bktable}
\par\addvspace{4pt}

\kgbsub{10.4}{付録の使い方}

この本の巻末には、2つの付録があります。付録Aは学習指導要領の側から、付録Bはシリーズの側から引く表です。お子さんが学校でいま習っている単元から章を探すときは付録Aを、シリーズの章が学校のどの学年の内容に当たるかを知りたいときは付録Bを使ってください。

\begin{bkfigure}
\begin{tikzpicture}[x=1mm, y=1mm]
  \node[kgb box, minimum width=40mm, minimum height=11mm] (s) at (20.5,0) {学習指導要領の項目\\{\mdseries（学年・科目・番号）}};
  \node[kgb box, minimum width=40mm, minimum height=11mm, fill=bkVol0!8] (c) at (119.5,0) {シリーズの章\\{\mdseries（全152章）}};
  \draw[kgb arr] ([yshift=2.2mm]s.east) -- node[kgb lab, above=0.8mm, font=\sffamily\bfseries\footnotesize] {付録A：項目から章を引く} ([yshift=2.2mm]c.west);
  \draw[kgb arr] ([yshift=-2.2mm]c.west) -- node[kgb lab, below=0.8mm, font=\sffamily\bfseries\footnotesize] {付録B：章から学年・科目を引く} ([yshift=-2.2mm]s.east);
\end{tikzpicture}
\bkfigcaption{2つの付録の引き方}
\end{bkfigure}
